\documentclass[11pt]{article}
% ===========================================================================
%  Shared preamble for the AP Calculus AB exam-preparation worksheets.
%  Input right after \documentclass. Files beginning with "_" are not compiled
%  on their own. Navy and gold, to match the course's lesson pages.
% ===========================================================================
\usepackage[letterpaper,top=2.3cm,bottom=2.2cm,left=1.7cm,right=1.7cm,headheight=28pt,footskip=30pt]{geometry}
\usepackage{amsmath,amssymb}
\usepackage[T1]{fontenc}
\usepackage[utf8]{inputenc}
\usepackage{lmodern}
\usepackage{xcolor}
\usepackage[most]{tcolorbox}
\usepackage{enumitem}
\usepackage{fancyhdr}
\usepackage{array}
\usepackage{tikz}
\usepackage{pgfplots}
\pgfplotsset{compat=1.18}

\definecolor{navy}{HTML}{161E54}
\definecolor{ink}{HTML}{27357E}
\definecolor{gold}{HTML}{B98A00}
\definecolor{goldbg}{HTML}{FFFDF3}
\definecolor{inkbg}{HTML}{F6F7FD}
\definecolor{rule}{HTML}{C9CFF0}
\definecolor{okgreen}{HTML}{1B7A44}
\definecolor{okbg}{HTML}{F1F8F3}

\pagestyle{fancy}
\fancyhf{}
\fancyhead[L]{\footnotesize Name: \underline{\hspace{4cm}}}
\fancyhead[C]{\textbf{\color{navy}Integration Academy}\\[-2pt]{\scriptsize Calculus AB \textperiodcentered{} AP\textsuperscript{\textregistered} Exam Preparation}}
\fancyfoot[L]{\tiny\color{black!55}\parbox[t]{0.78\textwidth}{AP\textsuperscript{\textregistered} is a registered trademark of the College Board, which was not involved in the production of, and does not endorse, this material. \textcopyright{} Integration Academy.}}
\fancyfoot[R]{\scriptsize\thepage}
\renewcommand{\headrulewidth}{1.2pt}
\renewcommand{\headrule}{\hbox to\headwidth{\color{gold}\leaders\hrule height \headrulewidth\hfill}}
\renewcommand{\footrulewidth}{0pt}
\setlength{\parindent}{0pt}
\setlength{\parskip}{4pt}

% Title block: \wstitle{code}{title}{unit line}
\newcommand{\wstitle}[3]{\begin{center}{\small\bfseries\color{gold}WORKSHEET #1}\\[3pt]{\LARGE\bfseries\color{navy}#2}\\[3pt]{\small\color{black!60}#3}\end{center}\vspace{2pt}}
% Section heading
\newcommand{\sech}[1]{\par\vspace{10pt}{\large\bfseries\color{navy}#1}\par\nobreak\vspace{2pt}{\color{rule}\hrule height0.8pt}\vspace{7pt}}
% A level tag printed beside a group of questions
\newcommand{\level}[1]{\par\vspace{6pt}{\footnotesize\bfseries\color{gold}\MakeUppercase{#1}}\par\nobreak\vspace{2pt}}

\newtcolorbox{ideabox}{colframe=ink,colback=inkbg,coltitle=white,colbacktitle=ink,fonttitle=\bfseries\small,title={Key ideas},arc=2pt,boxrule=0.8pt,left=8pt,right=8pt}
\newtcolorbox{example}[1]{colframe=rule,colback=white,coltitle=navy,colbacktitle=inkbg,fonttitle=\bfseries\small,title={#1},arc=2pt,boxrule=0.8pt,borderline west={3pt}{0pt}{ink},left=10pt,right=8pt}
\newtcolorbox{frbox}[1]{colframe=gold,colback=goldbg,coltitle=white,colbacktitle=gold,fonttitle=\bfseries\small,title={#1},arc=2pt,boxrule=0.8pt,breakable,left=8pt,right=8pt}
\newtcolorbox{solbox}{colframe=okgreen,colback=okbg,coltitle=white,colbacktitle=okgreen,fonttitle=\bfseries\small,title={Worked solutions},arc=2pt,boxrule=0.8pt,breakable,left=8pt,right=8pt}
\newtcolorbox{scorebox}[1]{colframe=navy,colback=white,coltitle=white,colbacktitle=navy,fonttitle=\bfseries\small,title={#1},arc=2pt,boxrule=0.8pt,breakable,left=8pt,right=8pt}

\newcommand{\soln}{\par\smallskip\textbf{\color{ink}Solution.}\ }
\newcommand{\result}[1]{\par\smallskip\textbf{#1}}
% blank working space of a given height
\newcommand{\work}[1]{\par\vspace{#1}}
% a point value in the margin of a scoring guide
\newcommand{\pts}[1]{\hfill{\footnotesize\bfseries\color{gold}[#1]}}

% Coordinate grid for graph questions: \begin{apgraph}{xmin}{xmax}{ymin}{ymax} ... \end{apgraph}
\newenvironment{apgraph}[4]{\begin{center}\begin{tikzpicture}\begin{axis}[
  axis lines=middle, xmin=#1, xmax=#2, ymin=#3, ymax=#4,
  grid=both, grid style={black!12}, xtick distance=1, ytick distance=1,
  ticklabel style={font=\scriptsize}, xlabel={$x$}, ylabel={$y$},
  width=10cm, height=7.2cm, samples=80, clip=false,
  every axis plot/.append style={line width=1.1pt, color=ink}]}{\end{axis}\end{tikzpicture}\end{center}}
\newcommand{\closeddot}[2]{\addplot[only marks, mark=*, mark size=2.2pt, color=ink] coordinates {(#1,#2)};}
\newcommand{\opendot}[2]{\addplot[only marks, mark=*, mark size=2.2pt, color=ink, mark options={fill=white}] coordinates {(#1,#2)};}

% Multiple-choice options: four across, or stacked when the options are long.
\newcommand{\opts}[4]{\par\vspace{3pt}\begin{tabular}{@{}p{0.235\linewidth}p{0.235\linewidth}p{0.235\linewidth}p{0.235\linewidth}@{}}(A)\ #1 & (B)\ #2 & (C)\ #3 & (D)\ #4\end{tabular}\par\vspace{9pt}}
\newcommand{\optsv}[4]{\par\vspace{3pt}\begin{tabular}{@{}l@{\ \ }p{0.9\linewidth}@{}}(A) & #1\\(B) & #2\\(C) & #3\\(D) & #4\end{tabular}\par\vspace{9pt}}

\fancyhead[R]{\footnotesize\color{black!60}Worksheet 4.6}
\begin{document}

\wstitle{4.6}{Tangent-Line Approximation}{Unit 4: Derivatives in Context}

\begin{ideabox}
\begin{itemize}[leftmargin=*,itemsep=2pt,topsep=2pt]
  \item Near $x=a$, \ $f(x)\approx L(x)=f(a)+f'(a)(x-a)$: the tangent line stands in for the curve.
  \item Choose $a$ where $f(a)$ and $f'(a)$ are easy and close to the value wanted.
  \item If the graph is \textbf{concave up} near $a$, the tangent line lies below it: an \textbf{underestimate}. Concave down: an \textbf{overestimate}.
  \item Justify over or under with the sign of $f''$ on the interval, never with the sign of $f'$.
\end{itemize}
\end{ideabox}

\sech{Worked examples}
\begin{example}{Example 1 \textperiodcentered{} A square root}
Use a tangent line to approximate $\sqrt{26}$.\soln $f(x)=\sqrt x$ at $a=25$: $f(25)=5$, $f'(25)=\dfrac{1}{10}$.\result{$\sqrt{26}\approx5+\dfrac{1}{10}(1)=5.1$.}
\end{example}

\begin{example}{Example 2 \textperiodcentered{} Over or under}
Use the tangent line to $f(x)=\sin x$ at $x=0$ to approximate $\sin(0.1)$. Is it an overestimate?\soln $L(x)=x$, so $\sin(0.1)\approx0.1$. $f''(x)=-\sin x<0$ for $0<x<0.1$.\result{$0.1$, an overestimate: the graph is concave down there.}
\end{example}

\begin{example}{Example 3 \textperiodcentered{} From values}
Given $f(2)=7$ and $f'(2)=-3$, approximate $f(2.2)$.\soln $7+(-3)(0.2)$.\result{$6.4$.}
\end{example}

\newpage
\sech{Practice}

\level{Core}
\begin{enumerate}[leftmargin=*,itemsep=2pt]
  \item Write the tangent-line approximation $L(x)$ for $f(x)=x^3$ at $x=2$.\work{2.3cm}
  \item Use your answer to Question 1 to approximate $(2.1)^3$.\work{2.3cm}
  \item Given $f(5)=4$ and $f'(5)=0.5$, approximate $f(4.8)$.\work{2.3cm}
  \item A graph is concave up near $x=a$. Is a tangent-line approximation there an overestimate or an underestimate?\work{2.3cm}
\end{enumerate}

\level{Exam level}
\begin{enumerate}[leftmargin=*,itemsep=2pt,start=5]
  \item Use the tangent line to $f(x)=\dfrac1x$ at $x=2$ to approximate $\dfrac{1}{2.1}$. Is it an overestimate or an underestimate? Justify.\work{2.8cm}
  \item Use the tangent line to $f(x)=\ln x$ at $x=1$ to approximate $\ln(0.9)$. Is it an overestimate or an underestimate? Justify.\work{2.8cm}
  \item A function has $f(1)=2$ and $f'(x)=x^2+\ln x$ for $x>0$. Approximate $f(1.1)$ and say whether the approximation is an overestimate or an underestimate.\work{2.8cm}
  \item The point $(3,1)$ lies on the curve $x^2+y^3=10$. Use the tangent line there to approximate the $y$-coordinate of the point on the curve with $x=3.1$.\work{2.8cm}
\end{enumerate}

\begin{frbox}{Question 9 \textperiodcentered{} Free-response practice}
Let $f(x)=\tan x$ for $-\dfrac\pi2<x<\dfrac\pi2$.
\begin{enumerate}[label=(\alph*),leftmargin=*,itemsep=2pt]
  \item Write an equation of the line tangent to the graph of $f$ at $x=\dfrac\pi4$.\work{2.6cm}
  \item Use the tangent line to approximate $\tan\!\left(\dfrac\pi4+0.05\right)$.\work{2.6cm}
  \item Is the approximation in part (b) an overestimate or an underestimate? Justify your answer.\work{2.6cm}
\end{enumerate}
\end{frbox}

\newpage
\begin{solbox}
\begin{enumerate}[leftmargin=*,itemsep=4pt]
  \item $L(x)=8+12(x-2)$.
  \item $8+12(0.1)=9.2$.
  \item $4+0.5(-0.2)=3.9$.
  \item An underestimate.
  \item $f(2)=\dfrac12$, $f'(2)=-\dfrac14$: $0.5-0.25(0.1)=0.475$. $f''(x)=\dfrac{2}{x^3}>0$, so concave up: an underestimate.
  \item $L(x)=x-1$, so $\ln(0.9)\approx-0.1$. $f''(x)=-\dfrac{1}{x^2}<0$, so concave down: an overestimate.
  \item $f'(1)=1$, so $f(1.1)\approx2.1$. $f''(x)=2x+\dfrac1x>0$ near 1: an underestimate.
  \item $2x+3y^2\dfrac{dy}{dx}=0$ gives $\dfrac{dy}{dx}=-\dfrac{6}{3}=-2$ at $(3,1)$, so $y\approx1-2(0.1)=0.8$.
  \item \begin{enumerate}[label=(\alph*),leftmargin=*,itemsep=3pt]
    \item $f\!\left(\dfrac\pi4\right)=1$ and $f'\!\left(\dfrac\pi4\right)=\sec^2\dfrac\pi4=2$: $y=1+2\left(x-\dfrac\pi4\right)$.
    \item $1+2(0.05)=1.1$.
    \item $f''(x)=2\sec^2x\tan x>0$ for $x$ near $\dfrac\pi4$, so the graph is concave up and the tangent line lies below it: an underestimate.
  \end{enumerate}
\end{enumerate}
\end{solbox}

\end{document}
